% ==============================================================================
% PHỤ LỤC B: CẨM NANG THÀNH PHẦN NHANH
%  -> Tra cứu khi viết; XÓA \input của file này trong main.tex trước khi nộp.
% ==============================================================================
\chapter{Cẩm nang thành phần nhanh}
\label{app:cam-nang}

Phụ lục này liệt kê các lệnh tiện ích của mẫu (định nghĩa trong
\texttt{config/commands.tex} và \texttt{config/colors.tex}). Mỗi khung gồm
\textbf{mã} (phần trên) và \textbf{kết quả} (phần dưới). Các gói nền tảng: tcolorbox
\cite{tcolorbox}; cú pháp LaTeX cơ bản xem \cite{lamport1994}.

\section{Chế độ cá nhân / nhóm}

Đổi một dòng ở đầu \texttt{main.tex}: \verb|\BaoCaoNhom| hoặc \verb|\BaoCaoCaNhan|.
Bản cá nhân dùng thành viên \textbf{đầu tiên} trong \texttt{metadata.tex}, ẩn bảng phân
công và đổi cách xưng hô.

\begin{vidulatex}[Từ ngữ tự đổi theo chế độ]
\ChungToi{} xin cam đoan... ; kết quả của \nhom{}.
Tùy ý: \NhomCaNhan{(bản nhóm)}{(bản cá nhân)}
\end{vidulatex}

\section{Nhấn mạnh trong đoạn văn}

\begin{vidulatex}
\nhanmanh{Ý chính}, \quantrong{rất quan trọng},
\danhdau{đánh dấu dạ quang}, \thuatngu{Học máy}[Machine Learning],
\nhan{Mới} \nhan[do]{Ưu tiên} \nhan[xanhla]{Hoàn thành},
\chuxanh{xanh} \chudo{đỏ} \chula{lá} \chucam{cam} \chutim{tím}.
\canbosung{số liệu năm 2025}
\end{vidulatex}

\begin{vidulatex}[Danh sách có biểu tượng]
\begin{dsdung} \item Đạt yêu cầu \end{dsdung}
\begin{dssai}  \item Chưa đạt     \end{dssai}
\begin{dssao}  \item Điểm nổi bật \end{dssao}
\end{vidulatex}

\section{Hộp thông báo và khung chữ}

\begin{vidulatex}
\begin{ghichu} Thông tin bổ sung. \end{ghichu}
\begin{luuy}[Lưu ý khi đọc số liệu] Nội dung. \end{luuy}
\begin{canhbao} Điều cần tránh. \end{canhbao}
\begin{goiy} Mẹo hoặc gợi ý. \end{goiy}
\end{vidulatex}

\begin{vidulatex}
\begin{ketqua} Kết quả/kết luận nổi bật. \end{ketqua}
\begin{dinhnghia}[Tên] Định nghĩa đánh số theo chương. \end{dinhnghia}
\begin{vidu} Ví dụ minh họa. \end{vidu}
\end{vidulatex}

\begin{vidulatex}[Khung tự chọn màu, khung nét đứt, trích dẫn]
\begin{khungchu}[tim]{Khung màu tím có tiêu đề} Nội dung. \end{khungchu}
\begin{khungchu}[ngoc] Khung không tiêu đề. \end{khungchu}
\begin{khungnetdut} Đề bài / yêu cầu. \end{khungnetdut}
\begin{trichdan}{Tác giả, năm} Câu trích dẫn. \end{trichdan}
\end{vidulatex}

\begin{vidulatex}[Ô chỉ số và so sánh hai cột]
\begin{hangchiso}[3]
  \ChiSo{95\%}{độ chính xác}
  \ChiSo[cam]{3,2s}{thời gian}
  \ChiSo[xanhla]{+12}{điểm}
\end{hangchiso}
\SoSanh{Ưu điểm}{Nhanh, rẻ.}{Nhược điểm}{Khó mở rộng.}
\SoSanh[xanhduong][tim]{Phương án A}{...}{Phương án B}{...}
\end{vidulatex}

\section{Công thức}

\begin{vidulatex}
\begin{khungcongthuc}[Công thức quan trọng]
  \begin{equation} \label{eq:vd}
    f(x) = \int_0^x e^{-t^2} \dif t
  \end{equation}
\end{khungcongthuc}
\begin{equation} \NoiBat{a^2 + b^2 = c^2} \end{equation}
Tham chiếu: \PT{eq:vd}; ký hiệu $\R, \N, \abs{x}, \norm{v}$.
\end{vidulatex}

\section{Bảng biểu}

\begin{vidulatex}[Bảng nhanh (bangbieu) — đổi thành bangsoc để có hàng sọc]
\begin{bangbieu}{Chú thích bảng}{tab:vd}{C{1cm} Y C{2.5cm}}
  \hangtieude \thead{STT} & \thead{Nội dung} & \thead{Giá trị} \\
  \midrule
  1 & Ô tốt  & \otot 9,5 \\
  2 & Ô xấu  & \oxau 2,0 \\
  3 & Ô nhấn & \onoibat 7,0 \\
\end{bangbieu}
\nguonbang{Tự tổng hợp.}
\end{vidulatex}

Kiểu cột có sẵn: \texttt{L\{rộng\}}, \texttt{C\{rộng\}}, \texttt{R\{rộng\}} (canh trái,
giữa, phải với độ rộng cố định), \texttt{Y}/\texttt{Z} (cột giãn đầy dòng, canh trái/giữa).
Bảng nhiều trang: xem \Bang{tab:bang-dai} ở \PhuLuc{app:du-lieu}.

\section{Hình ảnh và tham chiếu}

\begin{vidulatex}[Ảnh có khung; chưa có tệp -> hiện khung chờ]
\chenhinh[0.55\linewidth]{so-do-mau}{Chú thích hình}{fig:vd}
Tham chiếu: \Hinh{fig:vd}, \Bang{tab:vd}, \Chuong{chap:ket-qua},
\Muc{sec:muc-tieu}, \PhuLuc{app:du-lieu}.
\end{vidulatex}

Ảnh không khung: \verb|\chenhinh*[rộng]{tệp}{chú thích}{nhãn}|. Hai hình cạnh nhau
(tệp, chú thích, nhãn của từng hình rồi chú thích, nhãn chung):
\begin{lstlisting}[language={[LaTeX]TeX}, numbers=none]
\haihinh{a}{Chú thích a}{fig:a}{b}{Chú thích b}{fig:b}
        {Chú thích chung}{fig:ab}
\end{lstlisting}

\section{Mã nguồn và lệnh bash}

Khung mã: góc trái là \textbf{tên tệp}, góc phải là \textbf{badge ngôn ngữ}
(cũng là ngôn ngữ tô màu của \texttt{listings}: Python, C, C++, Java, JavaScript,
SQL, HTML, bash...). Tham số thứ ba (tùy chọn) là tùy chọn \texttt{listings}.
\begin{lstlisting}[language={[LaTeX]TeX}, numbers=none]
\begin{khungma}{main.cpp}{C++}
...
\end{khungma}
\end{lstlisting}
Kết quả:
\begin{khungma}{main.cpp}{C++}
#include <iostream>
int main() {
    std::cout << "Xin chào" << std::endl;   // in lời chào
    return 0;
}
\end{khungma}

\begin{khungma}{app.js}{JavaScript}[numbers=none]
const tong = (a, b) => a + b;   // không đánh số dòng
\end{khungma}

Lệnh bash dùng khung terminal riêng (dấu \texttt{\$} được tô xanh như dấu nhắc lệnh):
\begin{lstlisting}[language={[LaTeX]TeX}, numbers=none]
\begin{khungbash}[Tiêu đề cửa sổ]
...
\end{khungbash}
\end{lstlisting}
\begin{khungbash}[Biên dịch báo cáo]
$ cd HCMUT_02_Report_Class
$ latexmk -xelatex -outdir=build main.tex   # biên dịch bằng XeLaTeX
$ echo "Xong!"
\end{khungbash}

Đoạn mã có chú thích đánh số: dùng \texttt{lstlisting} với các tùy chọn
\texttt{language} và \texttt{caption} (xem \PhuLuc{app:du-lieu}).

\section{Lời giải}

\begin{vidulatex}[Khung lời giải: nền trắng xám, chữ xanh dương]
\begin{vidu} Tính $2 + 3 \times 4$. \end{vidu}
\begin{loigiai}
  Nhân trước, cộng sau: $2 + 3 \times 4 = 2 + 12 = 14$.
\end{loigiai}
\begin{loigiai}[Lời giải câu 2] Đổi tiêu đề bằng tham số tùy chọn. \end{loigiai}
\end{vidulatex}

\section{Bảng màu}

Mỗi màu có ba sắc độ: \texttt{<tên>}, \texttt{<tên>dam}, \texttt{<tên>nhat}.

\begin{center}
\setlength{\tabcolsep}{3pt}
\newcommand{\omau}[1]{\begin{tabular}{@{}c@{}}%
  \tikz{\fill[#1] (0,0) rectangle (0.9,0.45);
        \fill[#1dam] (0.9,0) rectangle (1.8,0.45);
        \fill[#1nhat] (1.8,0) rectangle (2.7,0.45);
        \draw[black!25] (0,0) rectangle (2.7,0.45);}\\[-2pt]
  \footnotesize\texttt{#1}\end{tabular}}
\omau{bkxanh}\omau{xanhduong}\omau{ngoc}\omau{xanhla}\omau{vang}\\[6pt]
\omau{cam}\omau{do}\omau{hong}\omau{tim}\omau{nau}
\end{center}
Màu riêng của trường: \texttt{bkxanh}, \texttt{bkxanhsang}, \texttt{bkxanhnen},
\texttt{bkvien}, \texttt{bkxam}. Thêm màu mới trong \texttt{config/colors.tex}:
\\ \verb|\definecolor{ten}{HTML}{RRGGBB}  \TaoSacDo{ten}|.
